

We define (see for instance Massimi \etal (2006) \cite{maqs06}) four reference 
quantities which are relevant for geodynamics\footnote{Note that in the paper 
the authors conflate $\rho$ and $\tilde{\rho}$ which prevents them from non-dimensionalising
all terms as we do here.}:
\begin{itemize}
\item a reference viscosity value $\underline{\eta}=10^{20}~\si{\pascal\second}$
\item a reference mass density $\underline{\rho}=1000~\si{\kg\per\cubic\metre}$
\item a reference time $\underline{t}=1\text{Myr}\simeq 3.15\cdot 10^{13}~\si{\second}$
\item a reference length $\underline{l}=1000~\si{\metre}$
\item a reference gravity $\underline{g}=9.81~\si{\metre\per\square\second}$
\end{itemize}
It follows that a reference pressure can be obtained:
\[
\underline{p}=\underline{\rho} \underline{g} \underline{l} = 9.81\cdot 10^6~\si{\pascal}
\]
Note that there is unfortunately no natural selection for the pressure scale. 
We could also have used $\underline{p}=\underline{\rho}\underline{\upnu}^2$ 
where dynamic effects are dominant i.e. high velocity flows,
or $\underline{p}=\underline{\eta}\underline{\upnu}/\underline{l}$ 
where viscous effects are dominant i.e. creeping flows (which 
is the case in geodynamics).
The definition of a reference velocity is more straightforward:
\[
\underline{\upnu} = \frac{\underline{l}}{\underline{t}} = 1~\si{\mm\per\year}
\]
We define dimensionless variables through:
\[
{\color{teal}x} = \frac{x}{\underline{l}}
\qquad
{\color{teal}y} = \frac{y}{\underline{l}}
\qquad
{\color{teal}z} = \frac{z}{\underline{l}}
\qquad
{\color{teal}\vec\upnu'} = \frac{\vec\upnu}{\underline{\upnu}}
\qquad
{\color{teal}t} = \frac{t}{\underline{t}}
\qquad
{\color{teal} \eta} = \frac{\eta}{\underline{\eta}}
\qquad
{\color{teal} g} =\frac{g}{\underline{g}}
\]
where the {\color{teal}teal} color indicates dimensionless values.

Consequently, time and space derivatives will be rescaled as follows:
\[
{\color{teal} \vec\nabla} = \underline{l}\; \vec\nabla
\qquad
{\color{teal} \partial_t} = \underline{t}\; \partial_t 
\]
Using this scaling relations the Navier-Stokes equation become:
\[
\frac{\rho \underline{l}}{\underline{t}^2} {\color{teal} \frac{\partial \vec\upnu}{\partial t}}
+
\frac{\rho \underline{l}}{\underline{t}^2} {\color{teal} (\vec\upnu \cdot\vec\nabla) \vec\upnu} 
=
- \underline{\rho} \underline{g} {\color{teal} \vec\nabla p}
+
\frac{\underline{\eta}}{\underline{l}\underline{t}} 
{\color{teal} \vec\nabla \cdot \eta (\vec\nabla \vec \upnu + \vec\nabla \vec \upnu ^T)}
+ \rho \vec{g}
\]
I make ${\color{teal}\rho}= \rho/\underline{\rho}$ appear in the left hand side:
\[
\frac{\underline{\rho} \underline{l}}{\underline{t}^2}
{\color{teal}\rho}
 {\color{teal} \frac{\partial \vec\upnu}{\partial t}}
+
\frac{\underline{\rho} \underline{l}}{\underline{t}^2} 
{\color{teal}\rho}
{\color{teal} (\vec\upnu \cdot\vec\nabla) \vec\upnu} 
=
- \underline{\rho} g {\color{teal} \vec\nabla p}
+
\frac{\underline{\eta}}{\underline{l}\underline{t}} 
{\color{teal} \vec\nabla \cdot \eta (\vec\nabla \vec \upnu + \vec\nabla \vec \upnu ^T)}
+ \rho \vec{g}
\]
which we can divide by $\underline{\rho} \underline{l}/\underline{t}^2$ to obtain:
\[
 {\color{teal}\rho} \left(
{\color{teal} \frac{\partial \vec\upnu}{\partial t}}
+
{\color{teal} (\vec\upnu \cdot\vec\nabla) \vec\upnu} 
\right)
=
- \frac{ \underline{g} \underline{t}^2 }{\underline{l}} {\color{teal} \vec\nabla p}
+
\frac{\underline{\eta} \underline{t}}{\underline{\rho} \underline{l}^2 } 
{\color{teal} \vec\nabla \cdot \eta (\vec\nabla \vec \upnu + \vec\nabla \vec \upnu ^T)}
+ \frac{\underline{t}^2}{ \underline{l}}  {\color{teal} \rho} \vec{g}
\]
One can recognise in this equation the Reynolds and Froude 
non-dimensional numbers (the ratio between the inertial and viscous forces, and the 
ratio between buoyancy and inertial forces respectively). 
\[
Re = \frac{\underline{\rho} \underline{l}^2}{\underline{\eta} \underline{t} }
\qquad
Fr= \frac{\underline{l}}{\underline{g}\underline{t}^2 }
\]
From this we conclude that inertial forces in the Earth's mantle
are small compared to viscous forces.
We can then write:
\[
\boxed{
{\color{teal}\rho} \left(
{\color{teal} \frac{\partial \vec\upnu}{\partial t}}
+
{\color{teal} (\vec\upnu \cdot\vec\nabla) \vec\upnu}  \right)
=
- \frac{1}{Fr} {\color{teal} \vec\nabla p}
+
\frac{1}{Re}
{\color{teal} \vec\nabla \cdot \eta (\vec\nabla \vec \upnu + \vec\nabla \vec \upnu ^T)}
+ \frac{1}{Fr} {\color{teal} \vec{g}}
}
\]
In our case, given the definitions taken above, we have:
\[
Re \simeq 3.174 \cdot 10^{-24}
\qquad
Fr \simeq 1.027 \cdot 10^{-25}
\]
so that the inertial terms can be dropped from the momentum equation (thereby yielding the 
dimensionless Stokes equations):
\[
{\color{teal} \vec\nabla \cdot \eta (\vec\nabla \vec \upnu + \vec\nabla \vec \upnu ^T)}
- \frac{Re}{Fr} {\color{teal} \vec\nabla p}
+ \frac{Re}{Fr}  {\color{teal} \vec{g} }
=0
\]
Note that in our case $Re/Fr\simeq 30.5$. 
